%=====================================================================
% Chapter: Experimental Setup
% Generated from the RecedingHorizon (RH-NBV) codebase configuration.
% Requires: amsmath, booktabs, siunitx (recommended), graphicx.
%=====================================================================
\chapter{Experimental Setup}
\label{ch:experimental-setup}

This chapter describes the hardware and software platform, the simulated
scene, the reconstruction and evaluation pipeline, the four viewpoint
planners under comparison, the measures taken to ensure a fair comparison
between them, and the family of occlusion scenarios used to probe planner
behaviour under progressively harder visibility conditions.

%---------------------------------------------------------------------
\section{Platform and Robot Setup}
\label{sec:platform}

All experiments were carried out on a 6-DoF \textbf{Universal Robots UR5e}
manipulator equipped with an \textbf{Intel RealSense} depth camera mounted
in an \emph{eye-in-hand} configuration, so that the camera pose is fully
determined by the arm's joint state. The system is implemented on
\textbf{ROS\,2 Jazzy}. Motion planning and execution are handled by
\textbf{MoveIt\,2}, and physics, rendering and depth sensing are simulated
in \textbf{Gazebo (Gz~Sim)}. The same planner code runs against the
physical robot; the simulator is used for the controlled comparative study
reported here because it provides a repeatable scene and exact ground truth.

The camera is modelled as a pinhole sensor with focal lengths
$f_x\approx 685.5$\,px and $f_y\approx 685.6$\,px and principal point
$(c_x,c_y)\approx(485.4,\,270.7)$\,px, producing registered colour and depth
images that are downsampled to a working resolution of $600\times450$\,px
for the reconstruction pipeline. The simulated depth sensor has a
\textbf{near clipping plane of \SI{0.40}{\meter}}: any surface closer than
this distance is removed from \emph{both} the depth and colour images. This
hardware limit drives the minimum stand-off constraint introduced in
Section~\ref{sec:fair}.

%---------------------------------------------------------------------
\section{Scene and Target Object}
\label{sec:scene}

The reconstruction target is the \textbf{Stanford Bunny} mesh, uniformly
scaled by a factor of $1.2$ and placed on a supporting surface in front of
the robot. In world coordinates its axis-aligned bounding box (AABB) spans
\[
  x\in[0.427,\,0.614],\quad
  y\in[-0.374,\,-0.229],\quad
  z\in[0.830,\,1.015]\ \text{(metres)},
\]
i.e.\ a footprint of roughly $18.7\times14.5$\,cm and a height of $18.5$\,cm.
The camera coordinate frame is oriented so that the optical (depth) axis
looks along the world $-Y$ direction toward the object.

A single \textbf{region of interest (ROI)} is defined as a cube of
half-width $\SI{0.095}{\meter}$ (\SI{\pm31}{voxels} at the working voxel
resolution) centred on the target position $(0.50,\,-0.30,\,0.92)$. All
coverage and reconstruction-quality metrics are computed strictly inside
this ROI, so that occluding structures placed just outside it (see
Section~\ref{sec:occlusion}) never inflate or pollute the scores.

%---------------------------------------------------------------------
\section{Scene Representation and Reconstruction Pipeline}
\label{sec:representation}

The scene is represented as a probabilistic \textbf{voxel grid}. The
reconstruction volume measures $0.3\times0.6\times0.3$\,m
($X\times Y\times Z$); it is deliberately elongated along $Y$ because that is
the camera depth axis. The voxel edge length is $\SI{3}{\milli\meter}$, and
the grid is centred on the ROI. At each planning iteration the segmented
depth image of the target is ray-cast into the grid, updating per-voxel
occupancy and yielding the observed voxel set used for scoring.

\paragraph{Coverage.}
ROI coverage is the fraction of ROI voxels that have been observed,
expressed as a percentage. Every planner begins each trial from an
identical, fully \emph{unobserved} ($0\%$) baseline.

\paragraph{Reconstruction quality (F1).}
Reconstruction fidelity is measured against the ground-truth mesh using
precision, recall and their harmonic mean, the \textbf{F1-score}, following
Burusa~et~al.~\cite{burusa2024}. A reconstructed occupied voxel is counted
as a true positive if it lies within a matching threshold of
$4\times$\,voxel size ($\SI{12}{\milli\meter}$) of the continuous mesh
surface; this tolerance accounts for the finite thickness of the occupancy
shell and is applied identically to all planners. When occluding panels are
present, voxels falling inside the (dilated) occluder bounding boxes are
excluded from the F1 computation for every planner alike, so that structure
belonging to the occluders is never rewarded or penalised as object
surface.

%---------------------------------------------------------------------
\section{Viewpoint Planners under Comparison}
\label{sec:planners}

Four planners are evaluated:

\begin{itemize}
  \item \textbf{RH-NBV} (\emph{proposed}) --- a receding-horizon,
        sampling-based next-best-view planner (Section~\ref{sec:rh}).
  \item \textbf{Gradient-NBV} --- the gradient-based information-driven
        baseline of Burusa~et~al.~\cite{burusa2024}.
  \item \textbf{PSO} --- a particle-swarm optimisation planner over the
        camera workspace.
  \item \textbf{Random} --- a naive baseline that samples admissible camera
        poses uniformly at random, providing a lower bound on performance.
\end{itemize}

\subsection{The Proposed Receding-Horizon Planner (RH-NBV)}
\label{sec:rh}

At each iteration the RH-NBV planner samples $K$ candidate \emph{sequences}
of $H$ future camera poses, scores each sequence by a discounted,
motion-penalised information-gain objective, and---in the spirit of
receding-horizon control---\emph{executes only the first pose} of the
best-scoring sequence before replanning from the newly observed state.

A candidate sequence
$\boldsymbol{\xi}=(\xi_1,\dots,\xi_H)$ is scored by
\begin{equation}
  J(\boldsymbol{\xi})
  \;=\;
  \sum_{k=1}^{H}
  \gamma^{\,k}
  \Bigl[\,
    f_{\mathrm{inc}}\!\left(M_{\mathrm{pred}},\,\xi_k\right)
    \;-\;
    \lambda\,C(\xi_{k-1},\xi_k)
  \,\Bigr],
  \label{eq:rh-objective}
\end{equation}
where $\gamma\in(0,1]$ is a discount factor that down-weights gains further
into the horizon, $f_{\mathrm{inc}}$ is the \emph{incremental}
information gain evaluated on a forward-predicted voxel grid
$M_{\mathrm{pred}}$ (counting only voxels not already accounted for by
earlier steps of the same sequence, which prevents the planner from being
rewarded for near-stationary sequences), $C(\cdot,\cdot)$ is the Euclidean
motion cost between consecutive poses, and $\lambda$ trades information gain
against travel. Candidate poses are drawn by mixing tangential ``orbit''
steps around the target (with probability equal to the bias ratio) with
uniform samples inside the admissible camera box; every sampled pose is
projected onto the admissible workspace and pushed out to the minimum
stand-off before evaluation. A stagnation-escape mechanism forces an
exploratory jump when coverage fails to improve for several consecutive
iterations. The default hyper-parameters are listed in
Table~\ref{tab:rh-params}.

\begin{table}[t]
  \centering
  \caption{Default RH-NBV hyper-parameters (all overridable per experiment).}
  \label{tab:rh-params}
  \begin{tabular}{@{}llc@{}}
    \toprule
    Symbol & Meaning & Value \\
    \midrule
    $H$        & Planning horizon (poses per sequence)      & $3$ \\
    $K$        & Candidate sequences per iteration          & $10$ \\
    $\gamma$   & Discount factor                            & $0.85$ \\
    $\lambda$  & Motion-cost weight                         & $2.0$ \\
    ---        & Step size (orbit tangential step, m)       & $0.065$ \\
    ---        & Bias ratio (orbit vs.\ uniform sampling)   & $0.70$ \\
    ---        & Occlusion-aware IG bonus                    & $2.0$ \\
    ---        & Stagnation patience (iterations)           & $4$ \\
    ---        & Stagnation threshold (\% coverage gain)    & $1.5$ \\
    \bottomrule
  \end{tabular}
\end{table}

%---------------------------------------------------------------------
\section{Ensuring a Fair Comparison}
\label{sec:fair}

Because the study's central claim concerns the planning \emph{strategy}, all
factors unrelated to strategy are held identical across the four planners
through a single shared configuration. Specifically, every planner uses the
same:

\begin{enumerate}
  \item \textbf{Reconstruction volume, ROI and F1 threshold} --- identical
        grid size, voxel resolution, ROI cube and $\SI{12}{\milli\meter}$
        matching tolerance, so the coverage denominator and quality metric
        are common to all.
  \item \textbf{Camera workspace} --- an axis-aligned box defined as the
        start pose $\pm(0.35,\,0.10,\,0.22)$\,m, with the $-Y$ face extended
        a further \SI{0.60}{\meter} toward and past the object. This
        ``wrap'' is essential: without it the camera box stays entirely on
        the $+Y$ (frontal) side of the object and no planner can reach the
        side views needed to see around occluders. The proposed planner's
        internal reach clamp is set \emph{equal} to this same box, so it can
        never search a smaller volume than the baselines.
  \item \textbf{Minimum stand-off} --- a hard lower bound of
        \SI{0.45}{\meter} on the camera-to-target distance, enforced for
        every planner. It sits just above the depth sensor's
        \SI{0.40}{\meter} near clip so that the target never vanishes from
        the depth and colour images.
  \item \textbf{Per-trial start perturbation} --- to model start/ROI
        uncertainty (as in~\cite{burusa2024}), trial $0$ starts from a
        fixed, deterministic pose while each subsequent trial $t$ applies a
        reproducible uniform perturbation of $\pm\SI{3}{\centi\meter}$ to
        the start position, seeded by \texttt{BASE\_SEED}$+t$
        (\texttt{BASE\_SEED}$=42$). Both the proposed planner and every
        baseline receive the \emph{same} perturbed start on a given trial,
        yielding a paired comparison with matched variance.
\end{enumerate}

Consequently the four planners search an identical physical camera space
under identical constraints; only their search \emph{strategy} differs.

%---------------------------------------------------------------------
\section{Experimental Protocol and Metrics}
\label{sec:protocol}

Each run consists of a fixed number of NBV iterations (\emph{viewpoints}).
Short runs use $5$ iterations and long convergence runs use $20$; the number
of independent trials per configuration is configurable, with results
averaged across trials. At every iteration the following quantities are
logged, and reported both per iteration (as curves) and as run-level
summaries:

\begin{itemize}
  \item \textbf{Primary:} ROI coverage (\%), F1-score, precision, recall.
  \item \textbf{Efficiency:} cumulative trajectory distance (m),
        ray-tracing calls (\#), and wall-clock computation time (s).
  \item \textbf{Aggregate / convergence:} coverage AUC (area under the
        coverage curve, rewarding \emph{how quickly} coverage is reached),
        coverage efficiency (\%/m), and views-to-threshold
        (iterations to reach F1$>0.7$ or coverage$>50\%$).
  \item \textbf{Reconstruction error:} Hausdorff and Chamfer distances
        between the reconstructed and ground-truth point sets.
  \item \textbf{Diagnostic:} stagnation count and first-detection iteration.
\end{itemize}

%---------------------------------------------------------------------
\section{Occlusion Scenarios}
\label{sec:occlusion}

To evaluate the planners under controlled and increasingly difficult
visibility conditions, the target is surrounded by rigid occluding
structures. Two families of scenarios are used.

\subsection{Legacy occluders}
A set of hand-authored worlds provides qualitatively distinct occlusion
types, selected at launch time via the \texttt{OCC} label:

\begin{description}
  \item[\texttt{none}] No occluder --- the unobstructed control condition.
  \item[\texttt{frontal}] A single panel directly between the object and the
        frontal ($+Y$) start region, forcing the planner to move to the
        sides.
  \item[\texttt{half\_box}] A partial box enclosing several faces of the
        object.
  \item[\texttt{tunnel}] Two opposing side walls forming a tunnel, leaving
        the front, back and top open.
  \item[\texttt{well}] A tight enclosure open only at the top, so the object
        is reachable mainly from above.
\end{description}

\subsection{Staged panel ladder}
\label{sec:panels}

The principal occlusion study uses a systematically generated
\textbf{staged panel ladder}, in which each stage closes exactly one
additional viewing ``aperture'' relative to the previous one, so that
difficulty rises \emph{monotonically} for every planner. The panels are
generated procedurally around the object's bounding box. All walls share an
identical thickness of $\SI{0.02}{\meter}$ and an identical height of
$\approx\SI{0.144}{\meter}$; each wall is placed with a $\SI{0.03}{\meter}$
clearance (gap) outside the object AABB, which keeps every wall
\emph{outside} the $\pm\SI{0.095}{\meter}$ evaluation ROI so that panel
voxels can never count toward object coverage. The walls are
bottom-anchored so that only the top $\approx\SI{5}{\centi\meter}$ of the
object (the bunny's head and ears) remains visible over the wall tops, and
wall widths span each face edge-to-edge so that adjacent walls \emph{meet at
the corners} with no gaps. The six stages are summarised in
Table~\ref{tab:panel-stages}.

\begin{table}[t]
  \centering
  \caption{The staged panel occlusion ladder. Each stage adds exactly one
    occluding element (bold), monotonically reducing visibility. Panels are
    \SI{0.02}{\meter} thick, $\approx\SI{0.144}{\meter}$ tall, offset
    \SI{0.03}{\meter} from the object AABB.}
  \label{tab:panel-stages}
  \begin{tabular}{@{}clp{7.2cm}@{}}
    \toprule
    Stage & Label & Occluding structure \\
    \midrule
    1 & \texttt{panels1} & One \textbf{side} wall only (front, back, other
        side and top all open). \\
    2 & \texttt{panels2} & \textbf{Tunnel}: the opposite side wall is added
        --- both sides closed, front/back/top open. \\
    3 & \texttt{panels3} & \textbf{U-shape}: the \textbf{front} wall is added
        --- the demanding frontal aperture closes. \\
    4 & \texttt{panels4} & \textbf{Four walls} (both sides, front, back);
        the top aperture remains open. \\
    5 & \texttt{panels5} & Four walls \textbf{+ bottom lid} flush under the
        walls; the top stays open. \\
    6 & \texttt{panels6} & \textbf{Sealed}: stage~5 plus a \textbf{top lid}
        resting on the wall tops (fully enclosed). \\
    \bottomrule
  \end{tabular}
\end{table}

The side from which the ladder starts is a difficulty knob. Starting from
the $+X$ (``right'') side is the \emph{easier} series, because the arm can
barely reach $+X$ viewpoints regardless; starting from the $-X$ (``left'')
side is \emph{harder}, because it blocks the side that is actually
arm-reachable. As a concrete example, at stage~4 the four generated walls
are centred at approximately $x=0.39$ and $x=0.65$ (side walls, spanning
$y$) and $y=-0.19$ and $y=-0.41$ (front and back walls, spanning $x$), all
at height $z\approx0.87$.

\paragraph{Consistent occluder masking.}
For every stage, the exact axis-aligned bounding boxes of the generated
panels (dilated by one voxel, \SI{3}{\milli\meter}) are written to a
per-stage manifest. All planners load this manifest and exclude the
corresponding voxels from F1 scoring, so the same regions are masked out for
every planner --- an essential condition for a fair comparison under
occlusion.
